\section{My team}

\textbf{2a. Structure, whether we discussed it, and how it changed.} \rtag{M16.b} Five people at kickoff, plus a late arrival. Myself (software/CV/autonomy); Robin (mission state-machine script, flight scripting on the Pi); Demetro (flight dynamics research, FDR presentation, aerodynamics write-up); Edward (hardware and CV calibration support); and our project manager / company pilot (risk assessments, field bookings, RC flying, logistics). In wk18 a sixth person, Zian, joined the group after their original team dissolved; Zian picked up mechanical tasks and CAD support, but because the software architecture was by then already concentrated on me, Zian's onboarding went through the mechanical side and he and I had very little direct collaboration -- a fact I list here because its absence is itself evidence of the problem \S3b diagnoses. At the wk12 kickoff we drew a rough ownership diagram on a whiteboard and each picked the work we felt competent to do. We did \textit{not} discuss interfaces, decision-making, contingency, or what happens when one person owns 70\% of the software surface area -- and that omission is the single biggest process lesson of the project. In Tuckman's terms \cite{tuckman1965} we skipped the \textit{storming} phase because nobody wanted to, and arrived at a \textit{norming} state that was really a polite silence. The real storming turned up in wk17 and wk22 instead, twice as hard.

Structure actually changed three times. First at \textbf{wk17}, when I presented a planned 11-state finite state machine (FSM) for the mission logic and two teammates (Robin and one other) pushed back hard that the architecture was too complex for anyone who had not written it to debug. I spent the evening re-reading their feedback, slept on it, and the next day refactored to a simpler FSM \cite{mainrefactor}. That episode -- the first time I genuinely changed a technical direction because of team input rather than my own evidence -- shifted my personal definition of ``good architecture'' from \textit{what I can debug} to \textit{what the team can debug}. I now see this is because I had been implicitly optimising for a reader who knew what I knew, which is the same cognitive bias McConnell flags in \textit{Code Complete} \cite{mcconnell2004} as the single largest reason code is ``unreadable to everyone except the author.'' I initially framed Robin's pushback as a \textit{preference} disagreement; in fact it was an \textit{evidence-of-failure} disagreement -- nobody but me could step through the FSM -- and I now see I preferred the preference framing because it protected me from admitting the code was badly adapted for its real users. Second, at \textbf{wk18} when Zian joined, the team re-allocated mechanical workstreams and I spent none of the resulting free time renegotiating my software footprint; in hindsight that was a missed structural opportunity. Third at \textbf{wk20}, after two cancelled flight days, we re-divided the remaining work: Robin would own the mission script end-to-end on the Pi, I would own passive observation, vision, GPS estimation, and a clean integration contract. I built \texttt{passive\_watch\_2.py} with a \texttt{-{}-robin} flag [commit \texttt{a5b1ab7}] and a JSON schema \cite{robinhandoff} specifically so Robin's state machine could poll vision over HTTP without touching the CV code [commits \texttt{885cee4}, \texttt{36b6069}, \texttt{24783d9}, \texttt{2bbadcec}, \texttt{56055e0}].

\textbf{2b. Most impactful team-working elements.} \rtag{M16.b} Ranked by how much they actually moved the needle, not how they felt at the time:

\begin{enumerate}
\item \textit{Explicit handoff contracts, written down.} When Robin and I agreed that his script owned flight and mine owned the camera, and I wrote the JSON schema and exact command-line invocation into \texttt{docs/DESIGN\_DECISIONS.md} and a README for his use, the integration happened without a single back-and-forth. That discipline did not exist in the first half of the project and arguments consumed hours; it started in wk20 and saved days. What looked like a technical problem for the first ten weeks was actually a trust-and-authority problem, because without a written contract every integration call became an implicit negotiation about whose mental model was authoritative.
\item \textit{Pairing with Edward on CV calibration.} Edward is strong on optics and hardware; I am strong on code. We agreed he would derive the field of view empirically (click-dummy at multiple altitudes in the DJI video, nine samples) and I would bake the result (54.4$^\circ$ horizontal FOV, focal 1416~px) into \texttt{config.py}. Neither of us could have produced that number in isolation in the time available. This was the only ``equal pairing'' I had on the project and its productivity-per-hour was noticeably higher than anything I did alone -- which is a data point I find hard to square with my own instinct to fall back on solo work whenever time is tight.
\item \textit{Pilot operating the ground station unaided on flight day 3.} Because \texttt{pi\_flight.py} ran in a browser with large, labelled buttons (Y/N/M/E/W/S) instead of keyboard shortcuts in a terminal, the pilot could actually operate it without me standing over his shoulder. That was the one moment on the project when my accessibility-focused design paid for itself with a human outcome rather than a technical one, and it lasted about forty minutes of flight-day uptime before a different hardware failure ended the session. The lesson persists even though the flight did not.
\item \textit{Psychological safety around the wk17 pushback.} Edmondson's work \cite{edmondson1999psychsafety} frames this better than I can: the reason Robin pushed back in the first place is that we had built just enough trust for a junior voice to challenge an architecturally-confident one. That is not an accident of personality; it is an outcome the team produced, and I now see it as the most important thing we did together, full stop.
\end{enumerate}

\textbf{2c. What I would change next time.} \rtag{M16.d-team} Three concrete things, each tied to a gap I now see clearly:

\begin{itemize}
\item \textbf{Contribution-design as a wk1 deliverable, not a wk20 retrofit.} Sit everyone down with the work-breakdown structure \textit{and} the interface list on day one, and decide who owns what \textit{contract}, not just what module. Measurable criterion: every teammate should be able to independently modify a module I wrote within one sprint. We did not hit this bar until wk20. The signal I will watch for next time is a teammate asking me ``can you just fix it?'' -- the moment I hear that question, the contract is wrong, not the person.
\item \textbf{Two pairs on software from wk1.} One person on one module is a single point of failure. I became that point of failure. Pairing would have slowed me individually by perhaps 30\%, but it would have bought the team a second person who could intervene when I was overloaded in wk22. Evaluating against velocity: the team's closed-task rate dipped from 9 to 4 tasks/week in the three weeks I was writing the complaint letter, which is a causal signature of exactly this failure mode.
\item \textbf{A shared persistent memory file (\texttt{MEMORY.md}-style) from wk1.} I maintained my own and it was transformative for my own productivity (blueprints, session archive, gotchas file). A shared version would have let teammates inherit my debugging history instead of re-learning it. Moon \cite{moon2004} argues reflection without a \textit{written} artefact tends to dissolve into impression; I now believe the same is true of team learning.
\end{itemize}

\textit{Conflict management, named explicitly.} The rubric rewards conflict management above a floor of 62 and I want to be specific rather than allude. Three episodes: (1) wk17, 11-state FSM vs simplified FSM -- Robin's pushback, my refactor overnight; resolved by my changing my mind rather than by negotiation, which is an incomplete resolution I come back to in \S3b. (2) wk19, ROS~2 vs raw MAVLink for inter-process communication -- I ``won'' by building a 20-line working heartbeat demo overnight ; that is persuasion-by-fait-accompli and I now label it a failure of leadership, not a success. (3) wk21, pilot workload -- the pilot told me the ground-station abort and confirm buttons were dangerously close; I separated them, added a confirmation dialog, and re-tested. The resolution mechanism in each case was different, and only case (3) involved me actually listening to the specific concern as stated, rather than substituting my own interpretation of it.
